% CHAPTER 5 - RESULTS AND DISCUSSION
\chapter{Results and Discussion}

\section{Experimental Setup}
All measurements in this chapter were taken on the development laptop described in
Section~3.9 (Intel Core i5-12450HX, 16\,GB RAM, NVIDIA GeForce RTX 3050 6\,GB laptop GPU,
Windows 11). The Python environment used Python 3.11.15, PyTorch 2.12 with CUDA available,
sentence-transformers 5.5.1, qdrant-client 1.12.1, sqlglot 25.2.0 and FastAPI 0.137.2.
These versions are newer than some pins in \texttt{requirements.txt} (for example
sentence-transformers 3.0.1 and FastAPI 0.111.0); the measured behaviour is that of the
installed versions. Because CUDA was available, sentence-transformers could place the
embedding and reranking models on the GPU.

Two kinds of evaluation were performed:
\begin{enumerate}[label=\arabic*., leftmargin=*, topsep=4pt, itemsep=3pt]
    \item The \textbf{automated test suite} (Section~4.8), which checks functional
          correctness, security rules and failure handling with mocked models.
    \item The \textbf{retrieval benchmark} \texttt{tests/run\_evaluation.py}, which uses the
          real \texttt{all-MiniLM-L6-v2} embedder and measures retrieval and context
          assembly on an in-memory platform database.
\end{enumerate}
The benchmark results stored in the repository (\texttt{tests/evaluation\_results.json})
were produced before the Qdrant vector store and the cross-encoder reranker were added.
For this report the benchmark was therefore run again, unchanged, on a separate copy of the
current code, where it used the embedded Qdrant store and the cross-encoder. Both results
are reported so that the effect of the retrieval upgrade is visible.

\section{Benchmark Design}

\subsection{Dataset}
The benchmark generates 120 synthetic cases in nine categories (Table~\ref{tab:dataset}).
All memories are stored for one user (\texttt{u-a1}) of organization \texttt{org-a} on one
connection, with a second user in the same organization and a user in a second organization
used for isolation checks. Memories accumulate over the run: each case adds its memory to
the same store, so later cases are retrieved against a growing set of distractors.

\begin{table}[H]
    \centering
    \small
    \caption{Composition of the 120-Case Benchmark}
    \begin{tabular}{p{4.6cm} p{0.9cm} p{7.9cm}}
        \toprule
        \textbf{Category} & \textbf{Cases} & \textbf{What is checked} \\
        \midrule
        Exact fact recall        & 15 & Question shares words with the stored definition. \\[3pt]
        Semantic fact recall     & 15 & Question paraphrases the definition (e.g.\ ``uncollectible receivables'' vs.\ ``write-offs''). \\[3pt]
        Business rule retrieval  & 15 & Procedural and semantic rules such as ``Display monetary results in INR''. \\[3pt]
        Cross-session recall     & 12 & KPI rules retrieved in a later, separate query. \\[3pt]
        Temporal recall          & 12 & After superseding, current query returns the new memory only and a history query returns the old one. \\[3pt]
        Conflicting memories     & 12 & After superseding, only the new definition is returned. \\[3pt]
        Follow-up understanding  & 15 & Prompt contains the previous analytical context for a follow-up. \\[3pt]
        Unrelated exclusion      & 12 & An unrelated HR memory is stored and an invoice question is asked. \\[3pt]
        User isolation           & 12 & A personal memory of user A1 must not be returned to user A2 or to another organization. \\
        \bottomrule
    \end{tabular}
    \label{tab:dataset}
\end{table}

\subsection{Metrics}
The metrics are computed exactly as implemented in the script. For the 57 retrieval cases
(exact, semantic, business rule and cross-session), a search with $k=3$ is made for the
case's query and the position of the target memory in the result is recorded:
\begin{align}
    \text{Recall@3} &= \frac{1}{N}\sum_{i=1}^{N} \mathbb{1}[\text{target}_i \in \text{top-3}_i] \\
    \text{MRR} &= \frac{1}{N}\sum_{i=1}^{N} \frac{1}{\text{rank}_i} \quad (0 \text{ if not retrieved})
\end{align}
following standard definitions \cite{manning2008ir, voorhees1999trec}.
\textbf{Precision@3} is implemented as $1/|\text{returned}|$ when the target is returned and
0 otherwise. With a single relevant memory per case and three results returned, its maximum
in practice is $1/3 \approx 0.33$, so low values are expected and should be read relative to
that ceiling. \textbf{Latency} is the wall-clock time of the search call.

\textbf{Context-inclusion accuracy} counts a case as correct when: the target memory text
appears in the built prompt (retrieval categories); superseding behaves as expected
(temporal and conflict categories); the prompt contains the previous analytical context and
question (follow-up category); or no memory is leaked (isolation category). In the
unrelated-exclusion category the script counts every case as correct without checking the
result, so these 12 cases do not measure anything. The \textbf{baseline} is a string
heuristic that marks a case correct only if it is an exact-fact case whose question already
contains ``excluding cancelled''; it represents a prompt without memory and is not an LLM
run. No LLM is called anywhere in the benchmark, so this metric measures whether the right
context reaches the prompt, \textbf{not} whether the generated SQL is correct.

\section{Retrieval Results}
Table~\ref{tab:retrieval} and Figure~\ref{fig:chart_retrieval} compare the stored run with
the re-run on the current pipeline.

\begin{table}[H]
    \centering
    \small
    \caption{Retrieval Metrics on the 57 Retrieval Cases}
    \begin{tabular}{p{5.2cm} p{3.6cm} p{3.6cm}}
        \toprule
        \textbf{Metric} & \textbf{Stored run} & \textbf{Re-run (current)} \\
        \midrule
        Recall@3                     & 0.7018 (40/57) & 0.7895 (45/57) \\[3pt]
        MRR                          & 0.6404         & 0.7281 \\[3pt]
        Precision@3 (as implemented) & 0.2485         & 0.2778 \\[3pt]
        Mean search latency          & 89.8\,ms       & 66.9\,ms \\
        \bottomrule
    \end{tabular}
    \label{tab:retrieval}
\end{table}

\begin{figure}[H]
    \centering
    \figimg{0.64\textwidth}{chart_retrieval.png}{Retrieval quality before and after adding Qdrant hybrid search and cross-encoder reranking.}{fig:chart_retrieval}
\end{figure}

With the current pipeline the target memory was among the top three results in 45 of 57
cases (Recall@3 = 0.79), compared with 40 of 57 before, and its average reciprocal rank
rose from 0.64 to 0.73. Precision@3 also increased, from 0.25 to 0.28, against a practical
ceiling of 0.33. Mean search latency was 67\,ms per query including embedding and reranking;
the lower value than the stored run should not be over-interpreted, because the two runs
were made at different times and the stored run's environment is not recorded.

\section{Context Assembly and Governance Results}
Table~\ref{tab:context_results} and Figure~\ref{fig:chart_inclusion} report the
context-level results.

\begin{table}[H]
    \centering
    \small
    \caption{Context-Inclusion, Update and Isolation Results}
    \begin{tabular}{p{6.4cm} p{3.0cm} p{3.0cm}}
        \toprule
        \textbf{Measure} & \textbf{Stored run} & \textbf{Re-run} \\
        \midrule
        Baseline context-inclusion (120 cases)      & 0.0083 (1/120) & 0.0083 (1/120) \\[3pt]
        CMF context-inclusion (120 cases)           & 0.9667 (116/120) & 0.9750 (117/120) \\[3pt]
        CMF, excluding the 12 unchecked cases       & 104/108 = 0.963 & 105/108 = 0.972 \\[3pt]
        Memory update accuracy (24 cases)           & 1.0 & 1.0 \\[3pt]
        Isolation pass rate (12 cases)              & 1.0 & 1.0 \\[3pt]
        Average prompt tokens, baseline             & 73.5 & 73.5 \\[3pt]
        Average prompt tokens, CMF                  & 613  & 747.9 \\
        \bottomrule
    \end{tabular}
    \label{tab:context_results}
\end{table}

\begin{figure}[H]
    \centering
    \figimg{0.78\textwidth}{chart_inclusion.png}{Context-inclusion accuracy and average prompt size, baseline versus CMF.}{fig:chart_inclusion}
\end{figure}

Three observations follow. First, the target business context reached the prompt in 105 of
the 108 checked cases with the current pipeline. Inclusion is higher than Recall@3 because
the prompt admits up to ten memories, not three. Second, all 24 superseding cases behaved
correctly: after an update, the current query returned the new definition and not the old
one, and history questions still found the old definition. Third, none of the 12 isolation
cases leaked a personal memory to another user or another organization; authorization
filters run before ranking, so a leak would require a filter bug rather than a ranking
error. The cost of richer context is prompt size: about 748 approximate tokens on average
against 74 for the schema-and-question baseline, which is well within the 4,000-token budget.

\section{Functional Test Results}
The complete suite was run with \texttt{pytest tests/}. All 64 collected tests passed in
7.92 seconds (Figure~\ref{fig:chart_tests}); the eight warnings reported include the
deprecation notice that the Gemini Python package used by the project is no
longer maintained. Together with the benchmark, the tests confirm the guardrail behaviour that the
benchmark does not cover: refusal of multi-statement, schema-changing, join-based and
unbounded changes, enforcement of grants, re-validation at approval time, masking of
sensitive columns for non-admins and read-only behaviour for Employees.

\begin{figure}[H]
    \centering
    \figimg{0.60\textwidth}{chart_tests.png}{Collected test cases per test file; all 64 passed.}{fig:chart_tests}
\end{figure}

\section{Qualitative Walkthrough}

\subsection{Governed Memories}
To illustrate the memory layer, three rules were entered on a copy of the application
connected to the retail sample database: two approved immediately (a revenue definition and
a currency preference) and one left pending (a definition of Gold-tier customers). The
Memories dashboard (Figure~\ref{fig:ui_memories}) shows the two approved memories with
their type, status, scope, visibility and subject, and counts one memory pending review. The
Evaluation page (Figure~\ref{fig:ui_evaluation}) displays the stored benchmark results and
the individual cases, and labels the accuracy figure as context inclusion rather than
executed SQL.

\begin{figure}[H]
    \centering
    \figimg{0.48\textwidth}{ui_memories.png}{Memories dashboard: two approved memories, one pending review.}{fig:ui_memories}\hfill
    \figimg{0.48\textwidth}{ui_evaluation.png}{Evaluation page with the stored benchmark results.}{fig:ui_evaluation}
\end{figure}

A memory search for ``monthly revenue by region'' then returned only the two approved
memories; the pending one was excluded as required. Table~\ref{tab:search_example} lists the
scores returned by the API. The revenue definition ranked first, helped by keyword overlap
(``revenue'') and a type match (``revenue'' maps to semantic memories); the currency
preference ranked second. In this example the cross-encoder assigned both short rule texts
a probability close to zero, so the final ranking was decided by the first-stage term. This
shows the value of keeping $(1-\lambda)S$ in the final score: the reranker refines the
order when it is confident, while the first-stage features still carry signal when it is
not.

\begin{table}[H]
    \centering
    \small
    \caption{Scores Returned for the Query ``monthly revenue by region''}
    \begin{tabular}{p{4.8cm} p{0.9cm} p{0.9cm} p{0.9cm} p{0.9cm} p{0.9cm} p{1.4cm}}
        \toprule
        \textbf{Memory} & $V$ & $K$ & $R$ & $I$ & $T$ & \textbf{Final} \\
        \midrule
        Revenue excludes cancelled and refunded orders \ldots & 0.344 & 0.333 & 1.0 & 0.778 & 1.0 & 0.2197 \\[3pt]
        Display monetary results in INR.                      & 0.349 & 0.000 & 1.0 & 0.778 & 0.4 & 0.1700 \\
        \bottomrule
    \end{tabular}
    \label{tab:search_example}
\end{table}

\subsection{Answer and Context Inspector}
Figure~\ref{fig:ui_chat_answer} shows a real answer produced earlier with Gemini on the
SaaS sample database to the question ``what is this data all about?''. The generated SQL
selects six descriptive columns from \texttt{accounts} with \texttt{LIMIT 5}; the answer
gives a direct summary, three insights, and states that it covers only the first five rows.
The sensitive \texttt{api\_key} column was not selected. The answer also shows that the
interface displays the model's Markdown markers literally, a presentation issue noted for
future work.

\begin{figure}[H]
    \centering
    \figimg{0.72\textwidth}{ui_chat_answer.png}{Plain-language answer with insights and the SQL that produced it.}{fig:ui_chat_answer}
\end{figure}

The Context Inspector for the same question (Figure~\ref{fig:ui_inspector}) shows that the
prompt used 1,170 of 4,000 tokens, contained no memories (none existed at that time) and one
recent turn, and begins with the instruction priority hierarchy and the security notice,
followed by the live schema with foreign-key references and sampled values.

\begin{figure}[H]
    \centering
    \figimg{0.72\textwidth}{ui_inspector.png}{Context Inspector: token usage, memories used and excluded, and the exact prompt.}{fig:ui_inspector}
\end{figure}

\section{Discussion and Threats to Validity}
The results support the design: authorization-first hybrid retrieval places the right
business context in the prompt in almost all checked cases, superseding keeps definitions
current, and isolation holds. Several limitations of the evidence must be stated plainly:
\begin{itemize}[leftmargin=*, topsep=4pt, itemsep=3pt]
    \item \textbf{No SQL accuracy was measured.} Whether the LLM applies an included rule
          correctly was not evaluated; context inclusion is a necessary, not a sufficient,
          condition for correct SQL.
    \item \textbf{Synthetic, template-generated data.} Several categories use near-identical
          templates (for example ``KPI\_$i$''), which favours keyword matching and may
          overstate performance on real, varied phrasing.
    \item \textbf{Weak baseline.} The baseline is a string heuristic, not a memory-less LLM
          run, so the difference between 0.8\% and 97.5\% shows that memories reach the
          prompt, not the size of an end-to-end improvement.
    \item \textbf{Unchecked cases and metric definitions.} The unrelated-exclusion category
          is counted as correct without a check, and Precision@3 is not the textbook
          definition; both were corrected for in the discussion above.
    \item \textbf{Single machine, single run.} Latency was measured once on one laptop.
\end{itemize}

\section{Performance Summary}

\begin{table}[H]
    \centering
    \small
    \caption{Summary of Measured Results (Current Pipeline)}
    \begin{tabular}{p{7.0cm} p{5.8cm}}
        \toprule
        \textbf{Measure} & \textbf{Result} \\
        \midrule
        Automated tests                         & 64 / 64 passed (7.92\,s) \\[3pt]
        Recall@3 / MRR (57 retrieval cases)     & 0.79 / 0.73 \\[3pt]
        Mean memory search latency              & 66.9\,ms \\[3pt]
        Context inclusion (108 checked cases)   & 97.2\% \\[3pt]
        Memory update accuracy (24 cases)       & 100\% \\[3pt]
        Isolation pass rate (12 cases)          & 100\% (no leaks) \\[3pt]
        Average prompt size, CMF vs.\ baseline  & 747.9 vs.\ 73.5 approx.\ tokens \\
        \bottomrule
    \end{tabular}
    \label{tab:summary}
\end{table}
